

\subsection{A PSPACE decision procedure for $\GWL$ 2}\label{app:pspace}


We sketch a non-deterministic PSPACE procedure to check validity in
$\GWL$. Let $\varphi$ be a formula. As discussed in
Sec.~\ref{sec_calculus}, if $\varphi\not\in\GWL$, there exists a
tree-like countermodel $\M=\stru{W,R,e}$ for $\varphi$ such that the
number of worlds is $O(\size{\varphi}^{\size{\varphi}})$.  

We can also bound the number $N$ of values needed to define $R$ and
$e$: in the worst case, the values $e(w,p)$ and $R(w,w')$ must all be
pairwise distinct, for every world $w$, $w'$ and every propositional
variable $p$ occurring in $\varphi$; accordingly, $N$ is in
$O(\size{\varphi}^{\size{\varphi}})$.  

Let $\Qcal=\{0\}\cup\{~1/k~|~1\leq k\leq N\}$; note that the elements
of $\Qcal$ can be non-deterministically generated in PSPACE.  By the
above discussion, if $\varphi\not\in\GWL$ there exists a countermodel
$\M$ such that the images of $R$ and $e$ are contained in $\Qcal$.  

To verify whether such a countermodel can be constructed, we introduce
a map $\Phi(w)$ that specifies the intended values of the subformulas
of $\varphi$ at world $w$.  More precisely, we assume that $\Phi(w)$
is data structure, representing a set of expressions of the form
$e(w,\a)=r$, where $w$ is a world name, $\a$ is a subformula of
$\varphi$, and $r\in\Qcal$.

We now sketch a nondeterministic algorithm for checking if a
countermodel for a formula $\varphi$ exists. During its execution, at
several points the algorithm selects a value from $\Qcal$, we call
such a selection a $\Wcal$-choice. Whenever it fails to find a
countermodel, it \emph{reconsiders the most recent open
  $\Qcal$-choice}, where \emph{open} means that not all
possible values in $\Qcal$ satisfying the requirements specified at
the point of the choice have yet been tried.

When the algorithm reconsiders an open $\Qcal$-choice:
\begin{enumerate}
\item a new value from $\Qcal$ satisfying the required conditions is
  selected;
\item the data structure $\Phi(w)$ is restored to the value it had at
  the time of the previous non-deterministic choice at that point.
\end{enumerate}
Note that to perform these task we only need ti remember at any
$\Qcal$-choice the next value to use (a rational value); for example
we can at any $\Qcal$-choice to select the next $n$ such that
$\frac{1}{n}\in\Qcal$ and $\frac{1}{n}$ satisfies the condition at the
$\Qcal$-choice point and to store it for the next time. Moreover, we
have to store the value of $\Phi(w)$ at the time of the
$\Qcal$-choice. Since that, given the world $w$, we have at most an
evaluation a $w$ for every subformula of $\varphi$, we have that every
$\Phi(w)$ we need to store ha length $O(|\varphi|)$ at most. The
algorithm cheks if a countermodel exists exploring the possible branches
of a countermodel generating at most $\size{\varphi}$ worlds. Since we
have at most a $\Qcal$-choice points for every possible subformula of
$\varphi$ at every possible worlds, we have that the number of
$\Qcal$-choices is bounden by $O(|\varphi|^2)$.  This means the maximum
amount of space required by the algorithm is $O(|\varphi|^3)$.


When, at some point of its execution, the algorithm fails and there is
no open non-deterministic choice to reconsider, it halts and we
conclude that the formula $\varphi$ has no countermodels, hence
$\varphi\in\GWL$. 

%Otherwise, the algorithm has generated during its
%execution, a sequence of worlds $\Phi(w_0),\dots,\Phi(w_k)$ which implicitly
%describe the evaluations if the  countermodel for $\varphi$.

\begin{enumerate}[label*=(Step~\arabic*)]
\item\label{step:1} The procedure considers the set
  $\Phi(w_0)=e(w_0,\varphi)=r_0$, where $w_0$ is the candidate root for a
  countermodel for $\varphi$ and $r_0$ is selected by a $\Qcal$-choice
  so that $r_0<1$. Then the procedure continues at~\ref{step:2}
  setting $\Phi(w)=\Phi(w_0)$.

\item\label{step:2} We apply propositional saturation steps to
  $\Phi(w)$ as long as possible. At each step, an expression of the
  form $e(w,\a\star\b)=r$ in $\Phi(w)$ is removed and replaced by
  evaluations of $\a$ and $\b$ that justify this value, where the
  corresponding values are selected by a $\Qcal$-choice.

  For instance, let $\psi=\a\to\b$.  
  \begin{itemize}
  \item If $r=1$, we must have $e(w,\a) \leq e(w,\b)$; we pick two
    values $r_1$ and $r_2$ with a $\Qcal$-choice so that $r_1\leq
    r_2$, and we replace $e(w,\a\to \b)=r$ in $\Phi(w)$ with $e(w,\a
    )= r_1$ and $e(w,\b)=r_2$.
  \item If $r<1$ we select a value $r_1$ with a $\Qcal$-choice so that
    $r_1> r$, and we replace $e(w,\a\to \b)=r$ in $\Phi(w)$ with
    $e(w,\a)= r_1$ and $e(w,\b) = r$.
  \end{itemize}
  At the end of this saturation phase, the formulas occurring $\Phi(w)$
  are either atomic or modal.
  
  \begin{itemize}
  \item If $\Phi(w)$ contains two evaluations $e(w,\a)=r_1$ and
    $e(w,\a)=r_2$ such that $r_1\neq r_2$, the construction
    \textsc{fails} and the algorithm must reconsider the most recent
    open $\Qcal$-choice, if any. If no open $\Qcal$-choice remains,
    the algorithm terminates and returns \textsc{fail}, meaning that
    no countermodel for $\varphi$ exists.

  \item If $\Phi(w)$ contains only atomic constraints, the
    construction succeeds and the algorithm returns \textsc{success}
    together with $\Phi(w)$.

  \item Otherwise, if $\Phi(w)$ contains at least the evaluation of a
    modal formula and we continue with~\ref{step:3}.
  \end{itemize}  

  % hence  $\varphi\not\in\GWL$.
\item\label{step:3} Let $e(w,\star\a_1)=r_1,\dots,e(w,\star\a_n)=r_n$
  the evaluations in $\Phi(w)$ involving modal formulas. 

  Select the next $k\in\{1,\dots,n\}$ and define $\Phi(w_k)$ as
  follows. Assume that $e(w,\star\a_k)=r_k$ is $e(w,\Box\a)=r$ (the
  case $e(w,\Diam\a)=r$ is similar). We generate a new world $w_k$ and
  select, by a $\Qcal$-choice, a value $r'_k$ for $R(w,w_k)$ and a
  value $r_k$ for $e(w_k,\a)$ such that $r'_k\to r_k = r$. The world
  $w_k$ is intended to represent the world witnessing that
  $e(w,\Box\a)=r$.

  We construct a new set $\Phi(w_k)$ by first adding the valuation
  ${e(w,\a_k)=r_k}$ and then extending it with the valuations required
  to justify the evaluations $e(w,\star'\b)=r'$ occurring in
  $\Phi(w_k)$. More precisely:

  We build up a new set $\Phi(w_k)$ and we add to $\Phi(w_k)$ the
  valuation $\{e(w,\a_k)=r_k\}$.  We then add to $\Phi(w_k)$ the
  valuations required to justify all the $e(w,\star'\b)=r'$ in
  $\Phi(w_k)$.  More precisely:
  \begin{itemize}
  \item for every $e(w,\Box \b)=r'$ in $\Phi(w)$, with $\b\neq \a$, we
    select with a $\Qcal$ choice a value $r''$ such that $r'_k \to r''
    \geq r'$ and we add $e(w_k,\b)=r''$ to $\Phi(w_k)$;
    

  \item for every $e(w,\Diam \b)=r'$ in $\Phi$, we select with a
    $\Qcal$-choice a value $r''$ such that $r'_k \land r'' \leq r'$ and
    we add $e(w_k,\b)=r''$ to $\Phi(w_k)$.
  \end{itemize}
  
  We now restart from~\ref{step:2} taking $\Phi(w)=\Phi(w_k)$. If the
  execution of~\ref{step:2} for $\Phi(w_k)$ \textsc{succeeds}, the
  algorithm proceeds to the next value of $k$, if any; if all values
  of $k$ have been considered, \ref{step:3} termites with
  \textsc{success}.

  Otherwise, if the execution of~\ref{step:2} for $\Phi(w_k)$
  \textsc{fails}, reconsider the most recent open $\Qcal$-choice, if
  any. If no open $\Qcal$-choice remains, the algorithm terminates
  with \textsc{fail}, meaning that no countermodel for $\varphi$
  exists.

  % Note that $\Phi(w_1)$ contains fewer modal operators than the input
  % formula $\varphi$; accordingly, we can recursively apply the
  % nondeterministic procedure to check whether a model having root
  % $w_1$ satisfying all the constraints in $\Phi(w_1)$ can be
  % constructed.  If the construction succeeds, we say that the
  % constraint $w_0:\Box\a=r$ is checked.
\end{enumerate}

% Note that, at each iteration, the information needed to check a constraint
% can be discarded, so the space usage remains polynomial in the size of $\varphi$.
% If all the modal constraints in $\Phi(w_0)$ are checked, we conclude that
% a countermodel for $\varphi$ can  be built, hence $\varphi\not\in\GWL$.
% If all the non-deterministic attempts to build a countermodel fail,
% we conclude $\varphi\in\GWL$.




% \newcommand{\ass}{\;\mapsfrom\;}

% %%% FUNCTION SATURATE
% \begin{function}[H]
%   {\sf Saturate}\Big($\Psi(w)$,\,$\Lambda(w)$\Big)\;
%   \tcp{param: $w:\a\star\b=r$ is a propositional constraint}
%   \tcp{param: $\Lambda$ is a list of valuataion maps for $w$}
%   $\Psi_0(w)\ass \Psi(w)$\;
%   $i\ass 0$\;
%   \While{($\Psi_i(w)$ contains a non atomic propositional constraint)}{ %%
%     Select a propositional non atomic constraint $w:\a\star\b=r$ in $\Psi(w)$\; %%
%     \eIf{($\exists r_1,r_2\in\Qcal$ new for $w:\a,w:\b$ 
%       w.r.t. $\Lambda(w)$ such that $r_1\star r_1=r$)} {%% THEN
%       $\Psi_{i+1}(w)\ass(\Psi_i(w)\setminus\{w:\a\star\b=r\})\cup\{w:\a = r_1,w:\b = r_2\}$\; %%
%     } { %% ELSE
%       \Return{FAIL}
%     }
%     $i\ass i+1$\;
%   }
%   \Return{$\Psi_i(w)$}
%   \caption{Saturate($\Psi(w)$,$\Lambda$)}\label{fun:saturate}
% \end{function}

% % %%% ModExp
% % \begin{function}[H]
% %   {\sf ModExp}\Big($\g$,\Phi(w)$\Big)\;
% %   %%\tcp{param: $w:\a\star\b=r$ is a propositional constraint}
% %   %%\tcp{param: $\Lambda$ is a list of valuataion maps for $w$}
% %   $\Psi_0(w)\ass \Psi(w)$\;
% %   % $i\ass 0$\;
% %   % \While{($\Psi_i(w)$ contains a non atomic propositional constraint)}{ %%
% %   %   Select a propositional non atomic constraint $w:\a\star\b=r$ in $\Psi(w)$\; %%
% %   %   \eIf{($\exists r_1,r_2\in\Qcal$ new for $w:\a,w:\b$ 
% %   %     w.r.t. $\Lambda(w)$ such that $r_1\star r_1=r$)} {%% THEN
% %   %     $\Psi_{i+1}(w)\ass(\Psi_i(w)\setminus\{w:\a\star\b=r\})\cup\{w:\a = r_1,w:\b = r_2\}$\; %%
% %   %   } { %% ELSE
% %   %     \Return{FAIL}
% %   %   }
% %   %   $i\ass i+1$\;
% %   % }
% %   % \Return{$\Psi_i(w)$}
% %   \caption{Saturate($\Psi(w)$,$\Lambda$)}\label{fun:saturate}
% % \end{function}


% %%FUNCTION MOD
% \begin{function}[H]
%   {\sf BuildMod}\Big($w:\varphi$,$\Psi(w)$\Big)\;
%   %\tcp{param: }
%   %\tcp{param: }
%   $i\ass 0$\;
%   $\Lambda_0(w)\ass []$\;
%   \While{($\exists r\in\Qcal$ new for $w:\a$ in $\Lambda_i(w)$)}{ %%
%     \eIf{($w:\varphi=r$ is a non atomic propositional constraint)}{ %% THEN
%       $\Psi(w)\ass ${\sf Saturate}($\Psi(w),\Lambda$)\;
%       \eIf{($\Psi(w)=FAIL$)}{%
%         $\Lambda_{i+1}\ass\Lambda_{i}\cup\{w:\varphi=r\}$
%       }{%
%         \eIf{($\Psi(w)$ contains a modal constraint $\gamma$)}{ %%
%           $\Psi(w)\ass BildMod(ModExp(\gamma,\Psi(w)))$
%           } { %% ELSE
%             \Return $\Psi(w)$
%           }%%
%         }
%     } {%% ELSE
%       $\Psi(w)\ass \{w:\varphi=r\}$\;
%     }
%   }
%   \caption{BuildMod}
% \end{function}


% \begin{itemize}
% \item A a pair of values $r_1,r_2$ in $\Qcal$ new w.r.t. $\Lambda(w)$,
%   if for every evaluation $\Psi(w)\Lambda(w)$ either $w:\a\not= r_1$ or $w:\b\not=r_2$.
% \end{itemize}


%%% Local Variables: 
%%% mode: latex
%%% TeX-master: "goedelModalLogicWitnessNonCrisp"
%%% End: 
